\documentclass[11pt]{article}
\usepackage[margin=1in]{geometry}
\usepackage{amsmath,amssymb,amsthm}
\usepackage[utf8]{inputenc}
\usepackage{listings}
\usepackage{hyperref}
\usepackage{xcolor}
\usepackage{textcomp}

\definecolor{leanblue}{rgb}{0.2,0.4,0.7}

\lstdefinelanguage{Lean}{
  keywords={theorem,lemma,def,axiom,structure,class,instance,where,by,sorry,simp,exact,intro,apply,have,let,fun,match,if,then,else,import,open,namespace,end,section,variable,inductive,noncomputable,abbrev,deriving},
  keywordstyle=\color{leanblue}\bfseries,
  comment=[l]{--},
  morecomment=[s]{/-}{-/},
  string=[b]",
  basicstyle=\ttfamily\small,
  breaklines=true,
  extendedchars=true,
  inputencoding=utf8,
  literate=
    {≥}{{$\geq$}}1
    {≤}{{$\leq$}}1
    {→}{{$\to$}}1
    {←}{{$\leftarrow$}}1
    {∧}{{$\wedge$}}1
    {∨}{{$\vee$}}1
    {¬}{{$\neg$}}1
    {∀}{{$\forall$}}1
    {∃}{{$\exists$}}1
    {⊆}{{$\subseteq$}}1
    {⊂}{{$\subset$}}1
    {∈}{{$\in$}}1
    {∉}{{$\notin$}}1
    {×}{{$\times$}}1
    {⊗}{{$\otimes$}}1
    {▸}{{$\blacktriangleright$}}1
    {⟨}{{$\langle$}}1
    {⟩}{{$\rangle$}}1
    {λ}{{$\lambda$}}1
    {α}{{$\alpha$}}1
    {β}{{$\beta$}}1
    {σ}{{$\sigma$}}1
    {θ}{{$\theta$}}1
    {ε}{{$\varepsilon$}}1
    {μ}{{$\mu$}}1
    {π}{{$\pi$}}1
    {Σ}{{$\Sigma$}}1
    {Π}{{$\Pi$}}1
    {▶}{{$\triangleright$}}1
    {⊔}{{$\sqcup$}}1
    {⊓}{{$\sqcap$}}1
    {⊥}{{$\bot$}}1
    {⊤}{{$\top$}}1
    {∅}{{$\emptyset$}}1
    {∪}{{$\cup$}}1
    {∩}{{$\cap$}}1
    {≠}{{$\neq$}}1
    {ℕ}{{$\mathbb{N}$}}1
    {₀}{{$_0$}}1
    {₁}{{$_1$}}1
    {₂}{{$_2$}}1
    {|>}{{$|{\!>}$}}2
    {·}{{$\cdot$}}1
    {—}{{---}}1
    {∘}{{$\circ$}}1
    {√}{{$\sqrt{}$}}1
    {∝}{{$\propto$}}1
    {≈}{{$\approx$}}1
    {═}{{=}}1
    {⟹}{{$\Longrightarrow$}}1,
}

\newtheorem{theorem}{Theorem}
\newtheorem{lemma}[theorem]{Lemma}
\newtheorem{definition}[theorem]{Definition}
\newtheorem{axiom_stmt}[theorem]{Axiom}

\title{Formal Verification: Verkle Tree State Commitments}
\author{Zach Kelling, Lux Industries}
\date{December 2025}

\begin{document}
\maketitle

\begin{abstract}
We formalize Verkle trees (vector commitment Merkle trees) used for Lux state commitments. Verkle trees use 256-way branching, reducing proof depth from $\log_2(n)$ to $\log_{256}(n)$. For $2^{32}$ accounts, depth drops from 32 to 4, yielding $5\times$ smaller proofs.
\end{abstract}

\section{Introduction}
Verkle trees replace Merkle trees for Lux state commitments. The 256-way branching factor dramatically reduces proof sizes, which is critical for light clients and cross-chain state verification. Each proof requires only $O(\text{depth})$ polynomial commitment elements.

\section{Theorems and Axioms}
\begin{axiom_stmt}[\texttt{verkle\_binding}]
Binding property: the same key cannot map to two different values under the same root.
\end{axiom_stmt}

\begin{theorem}[\texttt{depth\_for\_4b\_accounts}]
$256^4 = 4{,}294{,}967{,}296$: Verkle depth 4 supports $2^{32}$ accounts.
\end{theorem}

\begin{theorem}[\texttt{verkle\_proof\_savings}]
Proof size comparison: Merkle $1024$ bytes vs Verkle $192$ bytes yields $5\times$ savings.
\end{theorem}

\begin{theorem}[\texttt{branching\_advantage}]
Branching factor advantage: $256/2 = 128\times$ fewer levels per tree layer.
\end{theorem}

\begin{theorem}[\texttt{multiproof\_sublinear}]
Multiproofs grow sub-linearly: $k$ proofs share inner commitments.
\end{theorem}


\section{Lean Source}
\lstinputlisting[language=Lean]{../../formal/lean/Crypto/VerkleTree.lean}

\section{References}
\noindent Source code: \url{https://github.com/luxfi/proofs/blob/main/lean/Crypto/VerkleTree.lean}\\
\noindent White papers: \url{https://papers.lux.network}

\noindent Related white papers:
\begin{itemize}
  \item \texttt{lux-verkle-trees.tex}
\end{itemize}

\end{document}
